\documentclass[11pt,a4paper]{article}

% Packages
\usepackage[margin=2.5cm]{geometry}
\usepackage{amsmath,amssymb}
\usepackage{graphicx}
\usepackage{booktabs}
\usepackage{caption}
\usepackage{subcaption}
\usepackage{xcolor}
\usepackage{hyperref}
\usepackage{enumitem}
\usepackage{float}
\usepackage{multirow}

\hypersetup{
    colorlinks=true,
    linkcolor=blue,
    filecolor=magenta,
    urlcolor=cyan,
    citecolor=blue,
}

% Title info
\title{\textbf{Handwritten Digit Classification on MNIST Subset}\\
\large CS5487 Machine Learning -- Course Project Report}
\author{LAI Yikai\thanks{Student ID: 59563061, M.Sc. in Computer Science, City University of Hong Kong.} \and Zhang Yao\thanks{Student ID: 59843206, M.Sc. in Computer Science, City University of Hong Kong. Both students contributed equally to this project.}}
\date{April 2026}

\begin{document}

\maketitle

\begin{abstract}
This report presents an empirical study of multiple classification methods for handwritten digit recognition on a 4,000-sample MNIST subset. We compare eleven classifier configurations---including $k$-Nearest Neighbors ($k$-NN), Linear Discriminant Analysis (LDA), Quadratic Discriminant Analysis (QDA), Support Vector Machines (SVM) with linear and RBF kernels, Logistic Regression, and data-augmented variants---using 3-fold stratified cross-validation on the Trial~1 training data for model selection. Feature preprocessing techniques including scaling, Principal Component Analysis (PCA), and geometric data augmentation (rotation, translation, scaling) are investigated. The selected best model, SVM with RBF kernel combined with PCA to 35 components and mild data augmentation ($\times$3), achieves \textbf{96.37\%} mean accuracy on the MNIST test set (Trial~1: 96.50\%, Trial~2: 96.25\%), outperforming the 1-NN baseline of 91.60\%. On a held-out challenge dataset of 150 professor-written digits, the same model achieves \textbf{83.33\%}, a 15.03 percentage point improvement over the baseline. Detailed error analysis reveals that 1.6\% of samples are ``hard'' (misclassified by all models), while 14.0\% are ``medium'' where model capability differences matter.
\end{abstract}

\section{Introduction}

Handwritten digit recognition is a classic problem in machine learning and computer vision, serving as a benchmark for evaluating classification algorithms. The MNIST dataset \cite{lecun1998gradient} has been the standard testbed for decades. In this project, we work with a curated 4,000-sample subset of MNIST containing 400 images per digit class (0--9), where each image is a $28 \times 28$ grayscale pixel map vectorized into 784-dimensional feature vectors with pixel values in $[0, 255]$.

The project evaluates several classification methods covered in the CS5487 course, including:
\begin{itemize}[noitemsep]
    \item $k$-Nearest Neighbors ($k$-NN)
    \item Linear Discriminant Analysis (LDA)
    \item Quadratic Discriminant Analysis (QDA)
    \item Support Vector Machines (SVM) with linear and RBF kernels
    \item Logistic Regression
\end{itemize}

Additionally, we investigate the effect of feature preprocessing (normalization, PCA dimensionality reduction) and geometric data augmentation. We evaluate all trained models on a challenge dataset consisting of the professor's own handwritten digits, without any retraining or parameter tuning on the challenge data.

\section{Methodology}

\subsection{Dataset and Evaluation Protocol}

The dataset is stored in a MATLAB \texttt{.mat} file with the following fields:
\begin{itemize}[noitemsep]
    \item \texttt{digits\_vec}: $784 \times 4000$ matrix, each column is a flattened image.
    \item \texttt{digits\_labels}: $1 \times 4000$ vector of class labels $y \in \{0, \dots, 9\}$.
    \item \texttt{trainset}: $2 \times 2000$ matrix, two predefined training index sets.
    \item \texttt{testset}: $2 \times 2000$ matrix, two corresponding test index sets.
\end{itemize}

The dataset defines two fixed train/test splits (Trial~1 and Trial~2). In each trial, 2,000 samples (50\%) are used for training and 2,000 for testing. Importantly, the same writer does not appear in both training and test sets within a trial, ensuring writer-independent generalization. To prevent data leakage, model selection is performed exclusively on the Trial~1 training data via 3-fold stratified cross-validation. The selected model is then trained on the full training set and evaluated on the test set separately for each trial; the final reported accuracy is the mean of the two trial test accuracies.

\subsection{Classifiers}

\subsubsection{$k$-Nearest Neighbors ($k$-NN)}
$k$-NN is a non-parametric baseline classifier. We use Euclidean distance and test $k \in \{1, 3, 5\}$. No training is required; prediction is by majority vote among the $k$ closest training examples.

\subsubsection{Linear Discriminant Analysis (LDA)}
LDA assumes each class follows a Gaussian distribution with a shared covariance matrix. The decision boundary is linear. We use the standard eigenvalue decomposition solver.

\subsubsection{Quadratic Discriminant Analysis (QDA)}
QDA relaxes the shared covariance assumption, allowing each class its own covariance matrix. This yields quadratic decision boundaries. Since QDA requires estimating $10 \times 784 \times 784$ covariance parameters from only 2,000 training samples, we apply PCA to reduce dimensionality to 50 before fitting QDA. PCA is fit on the training data only.

\subsubsection{Support Vector Machines (SVM)}
SVM finds the maximum-margin hyperplane in a potentially high-dimensional feature space. We evaluate:
\begin{itemize}[noitemsep]
    \item \textbf{SVM-Linear}: Linear kernel, $C = 1.0$.
    \item \textbf{SVM-RBF}: Radial Basis Function kernel with $\gamma = \text{scale}$, $C = 1.0$.
    \item \textbf{SVM-RBF + PCA50}: RBF kernel after PCA to 50 components.
\end{itemize}
Multi-class SVM is implemented via a one-vs-all (OvA) wrapper: 10 binary classifiers are trained, one per digit. At test time, the classifier with the highest decision function value (furthest from the margin) is selected.

\subsubsection{Logistic Regression}
We use $L_2$-regularized multinomial logistic regression with $C = 1.0$ and a maximum of 1,000 iterations, also using the one-vs-all strategy. The features are standardized to zero mean and unit variance.

\subsubsection{Data Augmentation}
To improve generalization across different handwriting styles, we apply controlled geometric transformations to training images before normalization and PCA:
\begin{itemize}[noitemsep]
    \item \textbf{Rotation}: random angle in $[-5^\circ, +5^\circ]$ (mild), $[-10^\circ, +10^\circ]$ (medium), or $[-15^\circ, +15^\circ]$ (strong)
    \item \textbf{Translation}: random shift of $\pm 1$px (mild), $\pm 2$px (medium), or $\pm 3$px (strong)
    \item \textbf{Scaling}: random scale factor $0.95$--$1.05$ (mild), $0.90$--$1.10$ (medium), or $0.85$--$1.15$ (strong)
    \item \textbf{Strong mode}: additional Gaussian noise and brightness adjustment
\end{itemize}
Each training sample is augmented with $n$ synthetic copies, yielding $(1+n) \times 2{,}000$ training samples. Augmentation is applied \emph{before} normalization and PCA, and the same augmentation pipeline is used for both trials.

\subsection{Feature Preprocessing}

\begin{table}[H]
\centering
\caption{Preprocessing strategies per classifier.}
\label{tab:preprocessing}
\begin{tabular}{@{}llll@{}}
\toprule
\textbf{Classifier} & \textbf{Normalization} & \textbf{PCA} & \textbf{Augmentation} \\
\midrule
1-NN, 3-NN, 5-NN & None & None & None \\
LDA & None & None & None \\
QDA + PCA50 & None & 50 components & None \\
SVM-Linear & Scale to $[0, 1]$ & None & None \\
SVM-RBF & Scale to $[0, 1]$ & None & None \\
SVM-RBF + PCA50 & Scale to $[0, 1]$ & 50 components & None \\
Logistic Regression & Standard ($\mu=0, \sigma=1$) & None & None \\
SVM-RBF + PCA35 + Aug & Scale to $[0, 1]$ & 35 components & Mild $\times 3$ \\
SVM-RBF + PCA50 + Aug & Scale to $[0, 1]$ & 50 components & Mild $\times 3$ \\
\bottomrule
\end{tabular}
\end{table}

We experiment with three normalization schemes:
\begin{itemize}[noitemsep]
    \item \textbf{None}: Raw pixel values in $[0, 255]$.
    \item \textbf{Scale255}: Divide by 255 to map to $[0, 1]$.
    \item \textbf{Standard}: Per-feature standardization using training set statistics.
\end{itemize}
PCA is fit exclusively on the training data and applied to both training and test data to avoid data leakage.

\section{Experimental Setup}

All experiments were implemented in Python using scikit-learn \cite{pedregosa2011scikit}. The code is organized into modular components:
\begin{itemize}[noitemsep]
    \item \texttt{data\_loader.py}: Data loading, train/test splitting, normalization.
    \item \texttt{classifiers.py}: Classifier factory, one-vs-all wrapper, PCA transform.
    \item \texttt{augmentation.py}: Geometric data augmentation (rotation, translation, scaling).
    \item \texttt{experiments.py}: Experiment runner executing both trials.
    \item Analysis scripts: Error analysis, sample-level error analysis, class distribution visualization, and challenge evaluation.
\end{itemize}

\textbf{Evaluation Protocol.} We follow a two-phase evaluation protocol to avoid data leakage:
\begin{enumerate}[noitemsep]
    \item \textbf{Phase A (Model Selection):} All candidate methods are compared via 3-fold stratified cross-validation on the \emph{Trial~1 training data only} (2,000 samples). This includes comparing different classifiers, preprocessing pipelines, and hyperparameters. The configuration with the highest CV mean accuracy is selected.
    \item \textbf{Phase B (Final Evaluation):} The selected model is trained on the full training set of each trial and evaluated on the corresponding test set. The final reported accuracy is the mean of the two trial test accuracies with standard deviation.
\end{enumerate}

\textbf{Challenge Dataset.} A held-out challenge dataset of 150 professor-written digits is used to evaluate cross-writer generalization. The model selected in Phase~A is trained on each MNIST training split and evaluated on the challenge data \emph{without retraining or parameter tuning}.

Results are managed in a unified directory structure: \texttt{results/<datetime>/<script-name>/}, with a \texttt{meta.json} recording the git commit and timestamp, ensuring reproducibility.

\section{Experimental Results}

\subsection{MNIST Subset Classification Results}

Table~\ref{tab:cv-results} reports the 3-fold stratified cross-validation scores for all candidate methods on the Trial~1 training data. Table~\ref{tab:final-results} reports the final test accuracy of the selected model on both trials.

\begin{table}[H]
\centering
\caption{3-fold stratified cross-validation on Trial~1 training data (2,000 samples). Best CV score in \textbf{bold}.}
\label{tab:cv-results}
\begin{tabular}{@{}lcccc@{}}
\toprule
\textbf{Method} & \textbf{Fold 1} & \textbf{Fold 2} & \textbf{Fold 3} & \textbf{CV Mean $\pm$ Std} \\
\midrule
1-NN Baseline & 0.9175 & 0.9115 & 0.8904 & 0.9065 $\pm$ 0.0116 \\
3-NN & 0.9115 & 0.8966 & 0.8904 & 0.8995 $\pm$ 0.0089 \\
5-NN & 0.9130 & 0.8906 & 0.8934 & 0.8990 $\pm$ 0.0100 \\
LDA & 0.7271 & 0.7136 & 0.7252 & 0.7220 $\pm$ 0.0060 \\
QDA + PCA50 & 0.9265 & 0.9205 & 0.9189 & 0.9220 $\pm$ 0.0021 \\
SVM-Linear & 0.8441 & 0.8456 & 0.8483 & 0.8460 $\pm$ 0.0018 \\
SVM-RBF & 0.9355 & 0.9265 & 0.9174 & 0.9265 $\pm$ 0.0074 \\
SVM-RBF + PCA50 & 0.9520 & 0.9430 & 0.9324 & 0.9425 $\pm$ 0.0080 \\
Logistic Regression & 0.8396 & 0.8621 & 0.8634 & 0.8550 $\pm$ 0.0109 \\
\textbf{SVM-RBF + PCA35 + Aug} & \textbf{0.9670} & \textbf{0.9565} & \textbf{0.9520} & \textbf{0.9585 $\pm$ 0.0063} \\
SVM-RBF + PCA50 + Aug & 0.9670 & 0.9610 & 0.9474 & 0.9585 $\pm$ 0.0082 \\
\bottomrule
\end{tabular}
\end{table}

\begin{table}[H]
\centering
\caption{Final test accuracy of the selected model (SVM-RBF + PCA35 + Aug) on both predefined trials.}
\label{tab:final-results}
\begin{tabular}{@{}lccc@{}}
\toprule
\textbf{Method} & \textbf{Trial 1} & \textbf{Trial 2} & \textbf{Mean $\pm$ Std} \\
\midrule
\textbf{SVM-RBF + PCA35 + Aug} & \textbf{0.9650} & \textbf{0.9625} & \textbf{0.9637 $\pm$ 0.0012} \\
\bottomrule
\end{tabular}
\end{table}

\textbf{Key observations:}
\begin{enumerate}[noitemsep]
    \item \textbf{SVM-RBF + PCA35 + Aug} is selected by cross-validation (CV mean 95.85\%) and achieves 96.37\% mean test accuracy across both trials, outperforming the 1-NN baseline by 4.77 percentage points.
    \item Data augmentation provides substantial gains: both augmented SVM variants achieve the top CV scores ($\approx$95.9\%), while their non-augmented counterparts plateau around 94.3\%. Mild geometric transformations introduce valuable training diversity without distorting digit structure.
    \item PCA helps SVM-RBF significantly: CV accuracy improves from 92.65\% (no PCA) to 94.25\% (PCA50). PCA removes noisy high-frequency pixel variations, allowing the RBF kernel to focus on discriminative low-dimensional structure.
    \item QDA also benefits from PCA (92.20\% CV), as the full 784-dimensional covariance estimation is ill-posed with only 2,000 training samples.
    \item LDA performs poorly (72.20\% CV), suggesting the shared-covariance Gaussian assumption is too restrictive for this data.
    \item SVM-Linear (84.60\% CV) and Logistic Regression (85.50\% CV) are competitive but trail non-linear methods, indicating that digit classes are not linearly separable in pixel space.
    \item $k$-NN with $k > 1$ does not improve over 1-NN; in fact, 3-NN and 5-NN slightly underperform in CV, suggesting that the nearest neighbor alone is the most reliable predictor in this feature space.
\end{enumerate}

\subsection{Challenge Dataset Evaluation}

The challenge dataset contains 150 handwritten digits written by the course instructor. Models trained on the MNIST training split are evaluated directly on this data without retraining or parameter tuning. Table~\ref{tab:challenge-results} reports the results.

\begin{table}[H]
\centering
\caption{Challenge dataset accuracy (150 samples). Reference 1-NN baseline: 68.33\%.}
\label{tab:challenge-results}
\begin{tabular}{@{}lccc@{}}
\toprule
\textbf{Method} & \textbf{Trial 1} & \textbf{Trial 2} & \textbf{Mean $\pm$ Std} \\
\midrule
1-NN Baseline & 0.6600 & 0.7067 & 0.6833 $\pm$ 0.0233 \\
3-NN & 0.6667 & 0.6400 & 0.6533 $\pm$ 0.0133 \\
5-NN & 0.6333 & 0.6467 & 0.6400 $\pm$ 0.0067 \\
LDA & 0.4533 & 0.4200 & 0.4367 $\pm$ 0.0167 \\
QDA + PCA50 & 0.6867 & 0.7333 & 0.7100 $\pm$ 0.0233 \\
SVM-Linear & 0.6133 & 0.5533 & 0.5833 $\pm$ 0.0300 \\
SVM-RBF & 0.6867 & 0.7000 & 0.6933 $\pm$ 0.0067 \\
SVM-RBF + PCA50 & 0.7000 & 0.7333 & 0.7167 $\pm$ 0.0167 \\
Logistic Regression & 0.5800 & 0.6000 & 0.5900 $\pm$ 0.0100 \\
\textbf{SVM-RBF + PCA35 + Aug} & \textbf{0.8200} & \textbf{0.8467} & \textbf{0.8333 $\pm$ 0.0133} \\
SVM-RBF + PCA50 + Aug & 0.8200 & 0.8467 & 0.8333 $\pm$ 0.0133 \\
\bottomrule
\end{tabular}
\end{table}

\textbf{Key observations:}
\begin{enumerate}[noitemsep]
    \item \textbf{SVM-RBF + PCA35 + Aug} (the model selected by CV) achieves the best result (83.33\%), improving over the 1-NN baseline by 15.03 percentage points and over the non-augmented SVM-RBF + PCA50 by 11.67 percentage points.
    \item Data augmentation dramatically improves challenge generalization. While non-augmented models drop 20--25 percentage points from MNIST to challenge, the augmented model drops only 13.0 points (96.37\% to 83.33\%), demonstrating that geometric transformations help the model learn writer-invariant features.
    \item All non-augmented methods suffer a significant accuracy drop on the challenge data, indicating a substantial domain shift between student-written MNIST digits and professor-written challenge digits.
    \item QDA + PCA50 (71.00\%) is the best non-augmented method, confirming that PCA-based dimensionality reduction yields robust classifiers.
    \item $k$-NN degrades significantly on the challenge set, suggesting that the local neighborhood structure in pixel space does not transfer well across different writing styles.
\end{enumerate}

\subsection{Parameter Sensitivity Analysis}

To understand how hyperparameters affect performance, we conduct two controlled sensitivity studies using 3-fold stratified cross-validation on the training data only.

\subsubsection{Data Augmentation Intensity}

Table~\ref{tab:augmentation-cv} compares different augmentation configurations via cross-validation. The base model is PCA(35) + SVM-RBF($C=4$, $\gamma=\text{scale}$).

\begin{table}[H]
\centering
\caption{3-fold cross-validation scores for data augmentation strategies.}
\label{tab:augmentation-cv}
\begin{tabular}{@{}lcc@{}}
\toprule
\textbf{Configuration} & \textbf{CV Mean} & \textbf{CV Std} \\
\midrule
Original (no augmentation) & 95.05\% & 0.33\% \\
Mild augmentation $\times$3 & 95.85\% & 0.63\% \\
Medium augmentation $\times$3 & 96.50\% & 0.72\% \\
Strong augmentation $\times$3 & 96.35\% & 0.68\% \\
Medium augmentation $\times$5 & 96.45\% & 0.91\% \\
\textbf{Medium augmentation $\times$7} & \textbf{96.70\%} & \textbf{0.86\%} \\
\bottomrule
\end{tabular}
\end{table}

\textbf{Key findings:}
\begin{itemize}[noitemsep]
    \item \textbf{More augmentation generally helps.} Increasing augmentation intensity and copy count consistently improves CV scores, with medium augmentation $\times$7 achieving the highest score (96.70\%). Even strong augmentation (with added noise and brightness changes) reaches 96.35\%, outperforming the mild configuration selected in the main comparison.
    \item The selected model (mild $\times$3) was chosen from a predefined candidate set. These post-hoc sensitivity results suggest that stronger geometric transformations provide additional regularization benefits for this dataset.
    \item All augmented configurations substantially outperform the non-augmented baseline (95.05\%), confirming that geometric augmentation is beneficial regardless of intensity.
\end{itemize}

\subsubsection{SVM Hyperparameter and PCA Component Selection}

Table~\ref{tab:gridsearch} shows the result of a grid search over PCA components and SVM regularization parameter $C$, evaluated via 3-fold cross-validation on the training data.

\begin{table}[H]
\centering
\caption{GridSearchCV results (3-fold CV) for PCA components $\times$ SVM $C$. Best in \textbf{bold}.}
\label{tab:gridsearch}
\begin{tabular}{@{}lccccc@{}}
\toprule
\textbf{PCA} & $C=3$ & $C=4$ & $C=5$ & $C=6$ & $C=8$ \\
\midrule
20 & 95.20\% & 95.35\% & 95.25\% & 95.10\% & 95.15\% \\
25 & 96.00\% & 95.95\% & 95.80\% & 95.85\% & 95.60\% \\
30 & 96.00\% & 96.00\% & \textbf{96.25\%} & 96.15\% & 96.10\% \\
35 & 95.90\% & 95.85\% & 96.00\% & 96.10\% & 96.05\% \\
40 & 95.65\% & 95.55\% & 95.75\% & 95.70\% & 95.65\% \\
\bottomrule
\end{tabular}
\end{table}

\textbf{Key findings:}
\begin{itemize}[noitemsep]
    \item \textbf{Optimal region: PCA 30, $C=5$.} CV accuracy peaks at 96.25\%, with a broad plateau across PCA 25--35 and $C$ 3--6.
    \item Performance is relatively stable across the grid: the worst configuration (PCA=40, $C=4$ at 95.55\%) is only 0.70 percentage points below the best, indicating low sensitivity to these hyperparameters in the explored range.
    \item The chosen configuration (PCA=35, $C=4$) is near-optimal and generalizes well to both test trials, confirming that cross-validation on training data alone is sufficient for parameter selection.
\end{itemize}

\subsection{Error Analysis}

\subsubsection{Confusion Patterns}

We analyze the confusion matrices of SVM-RBF + PCA50 (the strongest non-augmented model analyzed) across both trials. The top confused digit pairs are:
\begin{itemize}[noitemsep]
    \item $5 \rightarrow 3$: 11 errors (2.75\% of digit 5 samples)
    \item $4 \rightarrow 9$: 9 errors (2.25\%)
    \item $9 \rightarrow 7$: 9 errors (2.25\%)
    \item $5 \rightarrow 6$: 8 errors (2.00\%)
    \item $9 \rightarrow 4$: 8 errors (2.00\%)
\end{itemize}

The most frequent confusions involve visually similar digits: 5 vs.~3 (both have a top curve and bottom loop), 4 vs.~9 (share an upper loop), and 9 vs.~7 (both have a top horizontal stroke).

Per-digit accuracy for SVM-RBF + PCA50 is high across all classes, with the lowest accuracy on digit 3 (93.00\%) and digit 9 (92.75\%). Digit 1 is the easiest (98.75\%), followed by digit 0 (97.75\%).

\subsubsection{Sample-Level Error Analysis}

To understand whether errors are sample-specific (hard samples) or model-specific, we compare predictions across five models (1-NN, 3-NN, SVM-RBF, SVM-RBF+PCA50, QDA+PCA50) on all 4,000 test samples:

\begin{itemize}[noitemsep]
    \item \textbf{Hard samples} (all 5 models wrong): 62 samples (1.6\%). These are likely genuinely ambiguous, poorly written, or potentially mislabeled digits. Model improvements will not help here.
    \item \textbf{Medium samples} (some models wrong): 560 samples (14.0\%). These indicate model capability differences---better models can improve on these.
    \item \textbf{Easy samples} (all models correct): 3,378 samples (84.5\%).
\end{itemize}

Among the medium samples, models agree on the wrong prediction for 460 samples but disagree for 100 samples, suggesting that for most errors, even the wrong predictions are structurally similar.

\subsection{Class Distribution}

The MNIST subset is perfectly balanced with 400 samples per digit class (10.0\% each). Both predefined train/test splits maintain this balance exactly (200 training / 200 test per class per trial). This eliminates class imbalance as a confounding factor in our experiments.

\section{Discussion}

\textbf{Effect of data augmentation.} Geometric augmentation is the single most impactful improvement in this study. The selected model (mild augmentation, $\times 3$) raises the CV score from 94.25\% (SVM-RBF + PCA50 without augmentation) to 95.85\% and the challenge accuracy from 71.67\% to 83.33\%. Post-hoc sensitivity analysis reveals that stronger augmentation configurations (medium $\times$3: 96.50\%, strong $\times$3: 96.35\%, medium $\times$7: 96.70\%) achieve even higher CV scores, suggesting that the predefined candidate set did not fully explore the optimal augmentation space. All augmented configurations substantially outperform the non-augmented baseline, confirming that geometric transformations help the model learn writer-invariant features.

\textbf{Effect of PCA.} PCA consistently improves non-linear methods (SVM-RBF and QDA). For the augmented SVM, reducing to 35 components is sufficient: the CV-selected model uses PCA35 and achieves performance comparable to PCA50. The 35--50 components explain approximately 80--83\% of the variance. By discarding high-frequency noise and redundancy in the 784-dimensional pixel space, PCA allows the classifiers to operate in a more structured, lower-dimensional manifold where class separation is cleaner.

\textbf{Effect of normalization.} Scaling pixel values to $[0, 1]$ is essential for SVMs with RBF kernels, as the kernel bandwidth $\gamma$ is computed relative to feature scales. Without scaling, pixel value magnitudes dominate the distance computation. Standardization (zero mean, unit variance) works well for Logistic Regression but is unnecessary for $k$-NN and tree-based methods.

\textbf{Linear vs.~non-linear.} The large gap between linear methods (LDA: 72.20\% CV, SVM-Linear: 84.60\% CV, Logistic Regression: 85.50\% CV) and non-linear methods (SVM-RBF: 92.65--95.85\% CV) confirms that digit classes have complex, non-linear decision boundaries in pixel space. This justifies the use of kernel methods and quadratic discriminant functions.

\textbf{Domain shift.} The challenge results highlight a critical issue in handwritten digit recognition: models trained on one writer population do not generalize perfectly to a different writer. Non-augmented models suffer a 20--25 percentage point drop in accuracy. However, data augmentation effectively closes this gap: the augmented model drops only 13.0 points (96.37\% to 83.33\%), demonstrating that geometric transformations help the model learn writer-invariant features. This suggests that data augmentation is a practical and effective technique for improving cross-writer generalization without requiring additional training data or domain adaptation methods.

\section{Conclusion}

We conducted a comprehensive empirical evaluation of eleven classifier configurations for handwritten digit classification, including data-augmented variants. Model selection was performed via 3-fold stratified cross-validation on the Trial~1 training data only, ensuring that Trial~2 remains an unbiased validation split. The challenge dataset was completely held out and used only for final evaluation. Our main findings are:
\begin{enumerate}
    \item \textbf{SVM with RBF kernel, PCA to 35 components, and mild data augmentation ($\times$3)} is selected by cross-validation (CV mean 95.85\%) and achieves 96.37\% mean test accuracy on the MNIST subset (Trial~1: 96.50\%, Trial~2: 96.25\%). On the challenge dataset, the same model achieves 83.33\%.
    \item \textbf{Data augmentation is the single most impactful improvement}: it raises challenge accuracy by 11.67 percentage points over the non-augmented model, effectively closing the generalization gap between student-written and professor-written digits. Post-hoc sensitivity analysis suggests that stronger augmentation configurations (medium and strong) can achieve even higher CV scores, indicating further optimization potential.
    \item \textbf{Hyperparameter sensitivity is low in the explored region.} GridSearchCV over PCA components [20, 25, 30, 35, 40] and SVM $C$ [3, 4, 5, 6, 8] shows a broad plateau with the peak at PCA~30 and $C=5$ (96.25\% CV). The selected configuration (PCA~35, $C=4$) is near-optimal, confirming robust parameter choices.
    \item PCA dimensionality reduction significantly benefits non-linear classifiers by removing noise and reducing the curse of dimensionality. Optimal component count is 35--50 (explaining ~80--83\% variance).
    \item Non-linear methods substantially outperform linear baselines, confirming that digit classes require flexible decision boundaries.
    \item Only 1.6\% of samples are universally misclassified (hard samples), while 14.0\% represent model-capability boundaries where better models can improve.
    \item The domain shift between MNIST and challenge data is substantial for non-augmented models (20--25 pp drop), but data augmentation reduces this gap to 13.0 pp, demonstrating its effectiveness for learning writer-invariant features.
\end{enumerate}

\section*{Contribution Statement}

Both group members contributed equally to this project. LAI Yikai was responsible for implementing the experiment pipeline, classifier wrappers, and the unified result management system. Zhang Yao contributed to the experimental design, error analysis scripts, and challenge evaluation. Both members participated in writing the project report and analyzing the results.

\begin{thebibliography}{9}
\bibitem{lecun1998gradient}
Y.~LeCun, L.~Bottou, Y.~Bengio, and P.~Haffner.
``Gradient-based learning applied to document recognition.''
\textit{Proceedings of the IEEE}, 86(11):2278--2324, 1998.

\bibitem{pedregosa2011scikit}
F.~Pedregosa et al.
``Scikit-learn: Machine learning in Python.''
\textit{Journal of Machine Learning Research}, 12:2825--2830, 2011.
\end{thebibliography}

\end{document}
